\section{Background: GPTSwarm}

\subsection{Agents as Graphs}

GPTSwarm defines an agent as a directed acyclic graph
\begin{equation}
    G=(N,E,F,o),
\end{equation}
where $N$ is a set of computational nodes, $E\subset N\times N$ is a directed edge set, $F=\{f_n\}_{n\in N}$ assigns a routine to each node, and $o\in N$ is an output node~\cite{pmlr-v235-zhuge24a}. A node consumes the original input and available predecessor outputs, applies its routine, and sends a natural-language or structured result to its successors. Topological execution makes edges both communication channels and control dependencies.

A swarm combines several agent graphs. Their internal edges remain fixed in the reported edge experiments, while additional edges between agents determine communication and orchestration. A task-specific decision node aggregates terminal information, for example by majority vote.

\subsection{Execution Architecture}

At execution time, the framework receives a task input, samples or fixes a feasible communication graph, and invokes nodes in topological order. Each node receives the original input plus messages from active predecessors. The decision node maps the terminal messages to a task output, after which an external evaluator produces utility. Graph optimization changes future edge probabilities; it does not differentiate through the language-model calls. This separation between graph sampling, black-box execution, and task evaluation is the architectural boundary retained by BC-GPTSwarm.

\subsection{Edge Optimization}

Let $\mathcal{E}=\{e_i\}_{i=1}^{d}$ be potential edges and let $D_{\theta}$ be a distribution over feasible DAGs. GPTSwarm optimizes
\begin{equation}
    \max_{\theta}\;\mathrm{E}_{G'\sim D_{\theta}}[u_{\tau}(G')],
\end{equation}
where $u_{\tau}$ is task utility. Each potential edge has a logit whose sigmoid is an inclusion probability. Edges are sampled sequentially and rejected when their addition would create a cycle.

Because LLM calls and task metrics are non-differentiable, GPTSwarm uses REINFORCE:
\begin{equation}
 \nabla_{\theta}\mathrm{E}[u_{\tau}(G)] \approx
 \frac{1}{M}\sum_{i=1}^{M}\hat{u}_{\tau}(G_i)
 \nabla_{\theta}\log p_{\theta}(G_i).
\end{equation}
The experiments use Adam with task-specific learning rates.

\subsection{Node Optimization}

For prompt-parameterized routines $f_n^{p_n}$, GPTSwarm accumulates node-level execution histories and applies an improvement operator to the current prompt, the operation description, and observed examples. The concrete HumanEval optimizer updates demonstrations and prompts using recent program-test outcomes. Node optimization is part of the original system but is outside the minimum Phase 0 reproduction scope.

\subsection{Published Evaluation}

The original study evaluates edge optimization on MMLU and Mini Crosswords, node optimization on HumanEval, and framework composition on GAIA. The appendix reports token, monetary, and wall-clock resource requirements, but these quantities do not enter the published optimization objective.
